\documentclass[12pt,a4paper]{ctexart}
\usepackage{geometry}\geometry{margin=2.2cm}
\usepackage{graphicx,float,fancyhdr,hyperref,longtable,booktabs,enumitem,xcolor}
\pagestyle{fancy}\fancyhf{}\fancyhead[L]{IO 服务器自主 OPC DA 采集技术方案}\fancyhead[R]{V1.0}\fancyfoot[C]{\thepage}
\title{\textbf{IO 服务器自主 OPC DA 采集\\技术方案}}
\author{约翰·刘}
\date{2026年7月13日}

\begin{document}\maketitle
\begin{abstract}
本文档阐述通过逆向工程力控 pSpace 6.0 实时数据库系统，实现对油田 IO 服务器（11.66.12.130:8889）的自主 OPC DA 数据采集方案。
核心技术路线包括：psAPISDK.dll 接口逆向、pSpace 协议分析、CommBridge 协议还原、12种保护装置设备建模，以及本机全栈模拟环境构建。
最终实现在不依赖原厂 IoProject/IoMonitor 的情况下，自主选择设备参数完成实时数据采集。
\end{abstract}

\section{项目背景}

\subsection{现状分析}
某油田 IO 采集系统基于力控 pSpace 6.0 平台构建，核心架构为：

\begin{verbatim}
  RTU(200+) → CommBridge(:53001) → IoMonitor → IoCommit → Oracle
  OPC DA → IOMan(TCP:8889) → pSpace(11.66.12.130) → IoCommit → Oracle
\end{verbatim}

系统已稳定运行多年，但存在以下问题：
\begin{enumerate}[leftmargin=*]
  \item 采集参数由 IoProject 固定分配，无法动态调整设备列表
  \item 依赖原厂 IoMonitor/IoCommit 全套组件，维护成本高
  \item pSpace 服务部署在远程服务器（11.66.12.130），开发调试不便
  \item 缺乏离线开发环境，新功能验证依赖生产网络
\end{enumerate}

\subsection{建设目标}
\begin{enumerate}[leftmargin=*]
  \item 实现自主选择设备参数完成 OPC DA 实时数据读取
  \item 构建本机全栈模拟环境，支持离线开发调测
  \item 建立完整的 IO 服务器知识本体，覆盖实体、关系、约束、闭环四层
  \item 提供五条可选的 OPC DA 数据访问路径
\end{enumerate}

\section{逆向工程成果}

\subsection{进程架构还原}
通过 WinRM 远程查询 131 生产服务器，提取了完整的进程快照：

\begin{table}[H]\centering
\caption{131 生产服务器进程清单}
\begin{tabular}{lrrl}
\toprule
\textbf{进程} & \textbf{PID} & \textbf{内存} & \textbf{功能} \\
\midrule
IoProject.exe & 5096 & 306MB & 总调度，进程编排 \\
IoMonitor.exe & 18400 & 61MB & GUI监视，Oracle写库 \\
CommBridge.exe & 19240 & 32MB & DTU/RTU桥接网关 \\
IOMan.exe ×8 & 多个 & -- & 设备采集器（OPC×7 + Modbus×1） \\
IoCommit.exe ×7 & 多个 & -- & 批量写Oracle \\
\bottomrule
\end{tabular}
\end{table}

\subsection{IOMan 命令行格式}
从 wmic 进程快照提取的 IOMan 启动参数：

\begin{verbatim}
IOMan.exe -aaa IO<SharedMemID>,<IoMonitorHwnd>,<DevType>,
                   <DBName>,<DevCount>:<DevID1>,<DevID2>,...

实例1 (Modbus): -aaa IO7CEE918,3B6C0B5C,1,IOCommitDB0,
                        10:02012170058,02105100097,...
实例2 (OPC):    -aaa IO7CEE898,3B6C0B5C,0,IOCommitDB0,
                        10:02110080028,02110110045,...
\end{verbatim}

\textbf{参数解析}（出自 IOMan.exe 格式字符串 \texttt{IO\%X,\%X,\%d,}）：
\begin{itemize}[leftmargin=*]
  \item SharedMemID：共享内存唯一ID（Global命名空间）
  \item IoMonitorHwnd：IoMonitor窗口句柄（WM\_COPYDATA IPC）
  \item DevType：0=OPC DA采集，1=Modbus RTU采集
  \item DBName：IOCommit数据库编号
  \item DevCount/DevIDs：设备数量及ID列表
\end{itemize}

\subsection{依赖链根因分析}
\begin{verbatim}
  pSpace Server (130:8889) ← 唯一远程依赖
      ↑ psAPI TCP协议
  IoMonitor.exe (可启动，不连pSpace则无功能)
      ↑ FindWindow + 共享内存
  IoProject.exe (无pSpace依赖！可独立运行)
      ↑ CreateProcess
  IOMan.exe (ConnectEx需要pSpace服务)
\end{verbatim}

\textbf{关键发现}：IoProject.exe 不依赖 psAPISDK.dll，CommBridge.exe 也不依赖。
这两个核心二进制可在本机独立运行。

\section{psAPISDK 接口逆向}

\subsection{SDK 来源}
从力控 pSpace 6.0.1.9 SDK 开发包获取完整头文件，包含：
\begin{itemize}[leftmargin=*]
  \item psAPIServer.h：服务器连接（Connect/Disconnect/IsConnected）
  \item psAPIReal.h：实时数据读写（ReadList/Subscribe/Write）
  \item psAPIHis.h：历史数据查询
  \item psAPIAlarm.h：告警管理
  \item psAPITag.h：标签管理和查询
  \item psAPIVariant.h：变体数据类型处理
  \item psAPICommon.h：基础类型定义（PSHANDLE=uint16，PSAPIStatus=int32）
\end{itemize}

\subsection{关键函数签名}
\begin{verbatim}
// 服务器连接
PSAPIStatus psAPI_Server_Connect(
    PSSTR pszServer,        // "IP:Port"，如"11.66.12.130:8889"
    PSSTR pszUserName,      // "admin"
    PSSTR pszPassword,      // "DQYTA11_pass"
    PSHANDLE *phServer);    // 返回连接句柄(uint16)

// 实时数据读取
PSAPIStatus psAPI_Real_ReadList(
    PSHANDLE hServer,
    PSUINT32 nCount,        // Tag 数量
    PSUINT32 *pTagIds,      // Tag ID 数组
    PS_VARIANT **ppValues,  // 返回值数组
    PSAPIStatus **ppErrors); // 错误码数组
\end{verbatim}

\subsection{PS\_DATA 内存布局（实测）}
\begin{table}[H]\centering
\caption{PS\_DATA 结构体内存布局（pack=4，24字节/记录）}
\begin{tabular}{lll}
\toprule
\textbf{偏移} & \textbf{大小} & \textbf{字段} \\
\midrule
0-3   & 4B & Time.Seconds (uint32 LE) \\
4-7   & 4B & Time.NanoSec (uint32 LE) \\
8-11  & 4B & DataType (uint32, =11) \\
12-15 & 4B & (reserved) \\
16-19 & 4B & \textbf{Float Value} (float32 LE) \\
20-23 & 4B & Quality (0xC0=GOOD) \\
\bottomrule
\end{tabular}
\end{table}

\section{五条 OPC DA 数据访问路径}

\begin{table}[H]\centering
\caption{OPC DA 数据访问路径对比}
\begin{tabular}{lp{3cm}ccp{2.5cm}}
\toprule
\textbf{\#} & \textbf{路径} & \textbf{实时} & \textbf{离线} & \textbf{状态} \\
\midrule
① & Python Mock 全栈模拟 & -- & 是 & 已验证(56设备) \\
② & dgiot\_lite Oracle 管线 & 否 & 否 & 已通(26万条) \\
③ & psAPISDK 直连 pSpace & 是 & 否 & \textbf{已通(自主读Tag)} \\
④ & IoProject 注入自定义 IOMan & 是 & 否 & 格式已验证 \\
⑤ & DCOM 直连 Kepware OPC DA & 是 & 否 & 脚本就绪 \\
\bottomrule
\end{tabular}
\end{table}

\subsection{路径③：psAPISDK 直连（核心方案）}
完整调用链：
\begin{verbatim}
  psAPI_Common_StartAPI()
  → psAPI_Server_Connect("11.66.12.130:8889", "admin",
                         "DQYTA11_pass", &handle)
    → handle=0x0001
  → psAPI_Real_ReadList(handle, 3, [5000,5001,5002],
                        &values, &errors)
    → values=[1.875, 1.875, 1.875]
    → timestamp=2026-07-13 19:35:06 (实时！)
    → quality=GOOD
\end{verbatim}

\textbf{实测数据}：
\begin{table}[H]\centering
\begin{tabular}{rrll}
\toprule
\textbf{Tag ID} & \textbf{值} & \textbf{时间戳} & \textbf{质量} \\
\midrule
5000 & 1.8750 & 2026-07-13 19:47:06 & GOOD \\
5501 & 2.2437 & 2026-07-13 19:46:59 & GOOD \\
6001 & 3.6946 & 2025-12-04 14:06:54 & 旧数据 \\
10000 & 3.0440 & 2026-07-11 20:12:20 & 旧数据 \\
\bottomrule
\end{tabular}
\end{table}

活跃 Tag ID 范围：5000--15000，值范围 1.7--3.7（工程值，A/kV/kW），
有效 Tag 以约 500 为间隔分布。

\subsection{路径①：本机全栈模拟}
\begin{verbatim}
  $ python tools/dev_env.py --scale 56

  服务           端口      状态
  ─────────────────────────────
  CommBridge     :53002    DTU桥接
  IoMonitor      :9002     数据汇聚
  IoCommit       :9003     数据写入
  IoProject      :9001     进程编排
  pSpace mock    :18889    pSpace数据源
  OPC DA mock    :13500    OPC数据源
  Modbus TCP     :502      直连轮询
\end{verbatim}

支持 20--100 台设备模拟，12 种设备类型全覆盖，带真实 Device.ini
通道配置和 ChangeData 转换系数。

\section{12种保护装置设备建模}
\begin{longtable}{lp{2.5cm}p{4cm}p{3cm}}
\toprule
\textbf{编码} & \textbf{名称} & \textbf{通道配置} & \textbf{转换系数} \\
\midrule
0x00 & DSL-31A 断路器 & 20通道 & C0=170/8192(电流/电压) \\
0x10 & DST-31A 变压器差动 & 15通道 & C1=8.5/8192(接地电流) \\
0x20 & DBPA-31A 备用电源 & 13通道 & C2=170/8192(相电压) \\
0x30 & DSB-31A 变压器后备 & 20通道 & C3=170×8.5/8192(功率) \\
0x40 & 电动机保护 & 19通道 & C4=1/8192(功率因数) \\
0x50 & DST-22D 变压器差动 & 20通道 & C5=2/8192(频率) \\
0x60 & DSB-22D 变压器后备 & 20通道 & C6-9=1(直通) \\
0x70 & DSL-24D 断路器 & 20通道 & -- \\
0x80 & DGP-11 变压器差动保护 & 21通道 & -- \\
0x90 & DGP-12 变压器后备保护 & 24通道 & -- \\
0xA0 & DGP-13 接地保护 & 22通道 & -- \\
0xB0 & DMP-31A 电动机保护 & 19通道 & -- \\
\bottomrule
\end{longtable}

物理量转换公式：$Y = raw \times Coefficient[i]$，基准标定值 8192(0x2000)。

\section{部署与使用}

\subsection{环境要求}
\begin{itemize}[leftmargin=*]
  \item 32位 Python 3.11（C:\textbackslash Python311-32\textbackslash）
  \item IO ServerOnLine 目录（含 psAPISDK.dll 及依赖 DLL）
  \item 网络可达 11.66.12.130:8889
\end{itemize}

\subsection{采集器使用}
\begin{verbatim}
# 读取指定 Tag ID
C:\Python311-32\python.exe tools/pspace_collector.py \
    --ids "5000,5501,6001,10000"

# 扫描活跃 Tag
C:\Python311-32\python.exe tools/pspace_collector.py \
    --scan --start 1 --end 20000 --step 100

# 启动离线模拟环境
python tools/dev_env.py --scale 20
\end{verbatim}

\section{风险与缓解}
\begin{table}[H]\centering
\begin{tabular}{lp{3cm}p{5cm}}
\toprule
\textbf{风险} & \textbf{概率} & \textbf{缓解措施} \\
\midrule
pSpace 密码变更 & 低 & 定期巡检连接状态 \\
130 网络中断 & 低 & 方案①离线模拟可替代 \\
Tag ID 映射缺失 & 高 & 待从 Oracle SYS\_POINTRELATION 提取 \\
32位 Python 兼容性 & 低 & 已封装独立采集器脚本 \\
\bottomrule
\end{tabular}
\end{table}

\section{总结}
本项目通过三周逆向工程，完整还原了力控 pSpace 6.0 系统的 IO 采集架构，
实现了不依赖原厂组件的自主 OPC DA 数据读取。
核心成果包括：
\textbf{psAPISDK 接口逆向}（3525 导出函数全量枚举）、
\textbf{pSpace 协议对接}（Connect→ReadList 全链路打通）、
\textbf{本机全栈模拟}（7 服务 + 56 设备离线可调）、
\textbf{五条数据访问路径}（涵盖 Mock/Oracle/psAPI/IOMan/DCOM）。

\end{document}
